% ============================================================
%  Chapter 11 – Further Techniques and Future Directions
% ============================================================

\section{Further Techniques and Future Directions}
\label{sec:future-directions}

The deep-zoom Mandelbrot rendering community has developed a rich
collection of techniques that complement or extend the algorithms
implemented in \FractalShark{}.  Understanding these methods is valuable
for two reasons: several of them directly influenced \FractalShark{}'s
design choices (e.g.\ the decision to use rebasing rather than
glitch correction, or linear approximation rather than series
approximation), and others represent promising avenues for future
improvement.

Each subsection below describes a technique, gives its key formulas, and
explains how it relates to---and differs from---\FractalShark{}'s current
approach.

% -----------------------------------------------------------
\subsection{Glitch Detection and Correction}
\label{subsec:glitch-detection}

Some deep-zoom renderers that lack per-pixel rebasing encounter
\emph{glitches}---pixels whose perturbation approximation has broken
down, producing visually incorrect results.  The standard remedy in
such renderers is a multi-pass detect-and-correct workflow.

The \emph{Pauldelbrot criterion} formalises glitch detection: a pixel is
considered glitched if, at some iteration~$n$,
\begin{equation}
  |z_n^\star| < |\Delta z_n|
  \label{eq:pauldelbrot-criterion}
\end{equation}
holds---that is, the magnitude of the reference orbit drops below the
perturbation delta, indicating that the reference is a poor fit.  The
correction procedure is to scan the rendered image for such pixels,
compute a new reference orbit centred on a glitched location, and
re-render only the affected region.  This process repeats until no
glitches remain and requires multi-reference orbit support
(\cref{subsec:multi-reference}).

\FractalShark{} sidesteps this workflow entirely by using
\emph{rebasing} (\cref{sec:perturb-only-rebase}): whenever
$|\Delta z_n| > |z_n|$ during iteration, the perturbation is reset by
folding the delta back into the reference and restarting the reference
index.  Rebasing corrects the instability per-pixel, on-the-fly, as
each pixel iterates---eliminating the need for post-render glitch
detection and multi-pass re-rendering.  In practice, rebasing is a more
direct and efficient solution: the problematic condition is handled
immediately at the point it occurs, rather than detected after the fact
and corrected in a separate pass.  Details of the glitch phenomenon and
related strategies are discussed in
\cite{HeilandAllen2013Perturbation} and \cite{HeilandAllen2022DeepZoom}.

% -----------------------------------------------------------
\subsection{Series Approximation (SA)}
\label{subsec:series-approx}

Series Approximation was the original technique for skipping iterations
in perturbation-based Mandelbrot rendering, introduced alongside
perturbation theory itself~\cite{Martin2013Perturbation}.  It has since
been largely superseded by linear approximation methods---both BLA
(\cref{sec:bilinear-approx}) and LA~v2
(\cref{sec:la-v2-perturb})---which are simpler, more robust, and faster
in practice.  SA is documented here for historical context.

SA constructs a polynomial $P(\Delta c)$ that approximates the
perturbation $\Delta z_N$ after $N$~iterations, so that each pixel can
skip those iterations with a single polynomial evaluation.  Starting
from the perturbation recurrence
\begin{equation}
  \Delta z_{n+1} = 2\,z_n^\star\,\Delta z_n + (\Delta z_n)^2 + \Delta c,
  \label{eq:sa-perturb-recurrence}
\end{equation}
one expands $\Delta z_n$ as a power series in $\Delta c$:
\[
  \Delta z_n = A_n\,\Delta c + B_n\,(\Delta c)^2
             + C_n\,(\Delta c)^3 + \cdots
\]
Substituting into \cref{eq:sa-perturb-recurrence} and matching powers of
$\Delta c$ yields recurrences for the coefficients:
\begin{align}
  A_{n+1} &= 2\,z_n^\star\,A_n + 1,
    \label{eq:sa-A} \\
  B_{n+1} &= 2\,z_n^\star\,B_n + A_n^2,
    \label{eq:sa-B} \\
  C_{n+1} &= 2\,z_n^\star\,C_n + 2\,A_n\,B_n,
    \label{eq:sa-C}
\end{align}
with higher-order terms following the same pattern.  Note that
\cref{eq:sa-A} is identical to the \code{CCoeff} recurrence used in
LA~v2 (\cref{sec:la-step-forward})---this is not a coincidence; both
capture the partial derivative $\partial z_n / \partial c$.

\paragraph{Why linear approximation superseded SA.}
SA and linear approximation (BLA, LA~v2) address the same
problem---skipping iterations that would otherwise be computed one at a
time---but they differ fundamentally in approach.  SA uses a
higher-order polynomial in $\Delta c$: a single evaluation can skip many
iterations, but the polynomial has no built-in validity bound and can
diverge catastrophically when the approximation radius is exceeded.
Renderers using SA must determine a safe skip count heuristically and
fall back to full perturbation if the polynomial produces nonsense.

Linear approximation, by contrast, uses first-order (linear)
coefficients with explicit, precomputed validity thresholds
(\cref{sec:la-threshold}).  BLA (\cref{sec:bilinear-approx}) builds a
generic power-of-two hierarchy of composable linear maps, agnostic of
the orbit's dynamics.  LA~v2 (\cref{sec:la-v2-perturb}) is more
sophisticated: it aligns the approximation hierarchy with the detected
period of the reference orbit, enabling each stage-0 entry to skip an
entire period and higher stages to skip multiples of the period.  Both
approaches degrade gracefully---when a validity threshold is
exceeded, the kernel falls back to perturbation for a few
iterations before attempting the next linear step.

In practice, the linear approach is simpler to implement (no
higher-order coefficient management), numerically more robust (bounded
error by construction), and often faster (the period-aligned hierarchy
in LA~v2 enables very large skip counts near periodic points without
risk of polynomial divergence).  Optimised SA variants such as NanoMB1,
NanoMB2, and Super Series Approximation (SSA)
(\cite{KallesFraktaler}) narrow the gap but remain more complex than
linear approximation for comparable or inferior performance.

% -----------------------------------------------------------
\subsection{Multi-Reference Orbit Strategies}
\label{subsec:multi-reference}

In renderers that use the glitch-detection workflow described in
\cref{subsec:glitch-detection}, a single reference orbit may not serve
all pixels well.  The multi-reference approach addresses this by
computing reference orbits at multiple locations and assigning each pixel
to the best available reference---the one that minimises $|\Delta z_n|$
at the iteration where the perturbation is most stressed.

Multi-reference rendering typically proceeds iteratively:
\begin{enumerate}
  \item Render the image with the primary reference orbit.
  \item Detect glitched pixels (\cref{subsec:glitch-detection}).
  \item Compute a new reference orbit centred on a glitched pixel's
        location.
  \item Re-render only the glitched pixels using the new reference.
  \item Repeat until no glitches remain.
\end{enumerate}

Because \FractalShark{} uses per-pixel rebasing
(\cref{sec:perturb-only-rebase}) rather than post-render glitch
detection, the motivation for multi-reference rendering is
substantially weaker: rebasing handles individual pixel instabilities
during iteration, avoiding the need for a separate detection-and-correction
pass.  Multi-reference strategies are primarily relevant to renderers
that lack rebasing and must instead rely on choosing a globally
well-suited reference for each region of the image.  The approach is
discussed in \cite{HeilandAllen2022DeepZoom}.

% -----------------------------------------------------------
\subsection{Smooth (Continuous) Iteration Count}
\label{subsec:smooth-coloring}

Standard escape-time coloring uses integer iteration counts, producing
visible banding in the rendered image.  The \emph{smooth iteration
count} eliminates this artifact by computing a fractional escape count.

For the standard bailout radius $R = 2$ the smooth count is
\begin{equation}
  n_{\text{smooth}} = n + 1
    - \frac{\log\!\bigl(\log|z_n|\bigr)}{\log 2},
  \label{eq:smooth-iter}
\end{equation}
where $n$ is the integer iteration at escape and $z_n$ is the final
orbit value.  The derivation rests on the potential function
$G(c) = \lim_{k\to\infty} \log|z_k| / 2^k$, which is continuous across
the exterior of the Mandelbrot set.  The escape time is related to the
level set $G(c) = \log R / 2^n$; interpolating between successive
iterates yields the fractional part.

For a general bailout radius~$R$ (note that \FractalShark{} uses
different squared-bailout thresholds in different contexts: $|z|^2 > 4$
in direct rendering kernels, $|z|^2 > 256$ in perturbation and
reference orbit paths, and $|z|^2 > 4096$ in the Feature Finder),
the formula generalises to
\begin{equation}
  n_{\text{smooth}} = n + 1
    - \frac{\log\!\bigl(\log|z_n|\bigr)}{\log 2}
    + \frac{\log\!\bigl(\log R\bigr)}{\log 2}.
  \label{eq:smooth-iter-general}
\end{equation}

This real-valued count can be used directly as a palette index,
eliminating the modular-integer banding visible in \FractalShark{}'s
current palette system (\cref{subsec:coloring}).  Implementation would
require modifying the antialiasing kernel (\cref{subsec:antialiasing})
to compute the fractional count from the final $|z_n|^2$ (already
available in the output buffer) and using it for palette lookup.  A
discussion of smooth coloring and other plotting algorithms appears in
\cite{WikipediaPlottingAlgorithms}.

% -----------------------------------------------------------
\subsection{Distance Estimation}
\label{subsec:distance-estimation}

The exterior distance estimator gives the approximate Euclidean distance
from a parameter~$c$ in the exterior of the Mandelbrot set to the
boundary $\partial\mathcal{M}$:
\begin{equation}
  d(c) \;\approx\; \frac{|z_n|\,\log|z_n|}{|z'_n|},
  \label{eq:distance-estimate}
\end{equation}
where $z'_n = \partial z_n / \partial c$ is the orbit derivative at the
escape iteration.

\FractalShark{} already computes $z'_n$: the perturbation pipeline
tracks $d_n = \partial z_n / \partial c$ for periodicity detection
(\cref{subsec:ref-orbit-periodicity}), and the LA~v2 kernel propagates
\code{dzdz} and \code{dzdc} through approximation steps
(\cref{sec:la-derivatives}).  The missing piece is the coloring
integration—mapping the distance $d(c)$ to a colour rather than (or in
addition to) the iteration count.

Applications of the distance estimator include:
\begin{itemize}
  \item Anti-aliased boundary rendering, where pixels near $\partial
        \mathcal{M}$ receive a smooth opacity falloff.
  \item Resolution-independent thin-filament rendering, ensuring that
        fine structures are visible regardless of pixel density.
  \item 3D height-field visualisation, using $d(c)$ as a height value
        to produce surface renders of the set.
\end{itemize}

Interior distance estimation is also possible—using the Koebe quarter
theorem and conformal mapping to the unit disk—but is substantially more
complex and requires detecting the attracting cycle first.  The
interaction of perturbation rendering with distance estimation is
discussed in \cite{HeilandAllen2013Perturbation}.

% -----------------------------------------------------------
\subsection{Interior Rendering}
\label{subsec:interior-rendering}

Points inside the Mandelbrot set—orbits that never escape—are
traditionally rendered as a uniform black.  These points belong to
\emph{hyperbolic components} (also called \emph{atom domains}), each
characterised by the period~$p$ of its attracting cycle.

Several interior coloring techniques exploit this structure:
\begin{itemize}
  \item \emph{Period-based coloring.}  Detect the period~$p$ of the
        attracting cycle and assign a colour according to~$p$.  This
        reveals the intricate structure of hyperbolic components.
  \item \emph{Phase coloring.}  Use the argument
        $\arg(z^\star)$ of the attracting fixed point~$z^\star$ to
        produce smooth colour gradients within a single component.
  \item \emph{Interior distance.}  Estimate how far the point lies from
        the boundary of its component, analogous to the exterior
        distance estimator (\cref{subsec:distance-estimation}).
\end{itemize}

Cycle detection during iteration can use Floyd's or Brent's
cycle-finding algorithms, which detect periodicity with $O(1)$ extra
memory.

\FractalShark{}'s Feature Finder (\cref{sec:feature-finder}) already
detects periods via Newton's method on the equation $z_p(c) = 0$; the
same period information could be used for interior coloring if the
detection logic were extended to run during rendering rather than as a
separate offline tool.

% -----------------------------------------------------------
\subsection{Boundary Tracing}
\label{subsec:boundary-tracing}

In a typical Mandelbrot image, large contiguous regions share the same
integer iteration count.  \emph{Boundary tracing} exploits this
observation: instead of computing every pixel independently, the
algorithm traces the contours between regions of different iteration
counts and flood-fills the interiors.

The procedure for each connected region of constant iteration count is:
\begin{enumerate}
  \item Identify a border pixel (one whose neighbour has a different
        iteration count).
  \item Follow the 4- or 8-connected border, recording the contour.
  \item Fill all enclosed pixels with the shared iteration count,
        skipping the per-pixel iteration entirely for interior pixels.
\end{enumerate}

The speedup depends on the fraction of the image occupied by large
uniform regions; typical gains range from $10\times$ to $50\times$ for
views dominated by the main cardioid or period-2 bulb.

\paragraph{Trade-offs.}
Boundary tracing is incompatible with smooth (continuous) coloring
(\cref{subsec:smooth-coloring}), because smooth coloring produces
continuously varying values with no discrete boundaries to trace.  The
technique works well with integer escape-time coloring.

Boundary tracing is orthogonal to perturbation and LA: it is a
pixel-level optimisation applied \emph{after} per-pixel iteration counts
are available for the border pixels; the per-pixel algorithm
(perturbation, LA, etc.)\ is unchanged.  Fraktaler-3
uses boundary tracing to achieve dramatic speedups.  An overview of the
algorithm is given in \cite{WikipediaPlottingAlgorithms}.

% -----------------------------------------------------------
\subsection{Histogram Equalization}
\label{subsec:histogram-equalization}

At deep zooms, iteration count distributions are highly skewed: most
pixels cluster around similar counts, with a long tail of rare high
values.  Standard palette cycling maps iteration counts to colours via
modular indexing, which wastes most of the colour palette on the
infrequent tail values and compresses the visually dominant range into a
narrow band.

\emph{Histogram equalization} addresses this by redistributing the
palette according to the actual distribution of iteration counts:
\begin{enumerate}
  \item Collect a histogram $H(k)$ of iteration counts across all
        pixels.
  \item Compute the cumulative distribution function
        $\mathrm{CDF}(k) = \sum_{j \le k} H(j) \big/ \sum_j H(j)$.
  \item Map each pixel's iteration count~$k$ to the colour at palette
        index $\mathrm{CDF}(k)$.
\end{enumerate}
The effect is that all colours in the palette are utilised approximately
equally, maximising visual contrast in the regions where most detail
exists.

\FractalShark{} already computes per-frame iteration statistics
(minimum, maximum, sum) via the reduction kernel
(\cref{subsec:reduction}); extending this to a full histogram would
require a histogram kernel but presents no architectural obstacles.  The
technique is independent of the iteration algorithm and can be combined
with smooth coloring (\cref{subsec:smooth-coloring}) for even better
results.  A general discussion of histogram-based coloring appears in
\cite{WikipediaPlottingAlgorithms}.
